% Options for packages loaded elsewhere
\PassOptionsToPackage{unicode}{hyperref}
\PassOptionsToPackage{hyphens}{url}
\documentclass[
  ignorenonframetext,
]{beamer}
\newif\ifbibliography
\usepackage{pgfpages}
\setbeamertemplate{caption}[numbered]
\setbeamertemplate{caption label separator}{: }
\setbeamercolor{caption name}{fg=normal text.fg}
\beamertemplatenavigationsymbolsempty
% remove section numbering
\setbeamertemplate{part page}{
  \centering
  \begin{beamercolorbox}[sep=16pt,center]{part title}
    \usebeamerfont{part title}\insertpart\par
  \end{beamercolorbox}
}
\setbeamertemplate{section page}{
  \centering
  \begin{beamercolorbox}[sep=12pt,center]{section title}
    \usebeamerfont{section title}\insertsection\par
  \end{beamercolorbox}
}
\setbeamertemplate{subsection page}{
  \centering
  \begin{beamercolorbox}[sep=8pt,center]{subsection title}
    \usebeamerfont{subsection title}\insertsubsection\par
  \end{beamercolorbox}
}
% Prevent slide breaks in the middle of a paragraph
\widowpenalties 1 10000
\raggedbottom
\AtBeginPart{
  \frame{\partpage}
}
\AtBeginSection{
  \ifbibliography
  \else
    \frame{\sectionpage}
  \fi
}
\AtBeginSubsection{
  \frame{\subsectionpage}
}
\usepackage{iftex}
\ifPDFTeX
  \usepackage[T1]{fontenc}
  \usepackage[utf8]{inputenc}
  \usepackage{textcomp} % provide euro and other symbols
\else % if luatex or xetex
  \usepackage{unicode-math} % this also loads fontspec
  \defaultfontfeatures{Scale=MatchLowercase}
  \defaultfontfeatures[\rmfamily]{Ligatures=TeX,Scale=1}
\fi
\usepackage{lmodern}
\ifPDFTeX\else
  % xetex/luatex font selection
\fi
% Use upquote if available, for straight quotes in verbatim environments
\IfFileExists{upquote.sty}{\usepackage{upquote}}{}
\IfFileExists{microtype.sty}{% use microtype if available
  \usepackage[]{microtype}
  \UseMicrotypeSet[protrusion]{basicmath} % disable protrusion for tt fonts
}{}
\makeatletter
\@ifundefined{KOMAClassName}{% if non-KOMA class
  \IfFileExists{parskip.sty}{%
    \usepackage{parskip}
  }{% else
    \setlength{\parindent}{0pt}
    \setlength{\parskip}{6pt plus 2pt minus 1pt}}
}{% if KOMA class
  \KOMAoptions{parskip=half}}
\makeatother
\setlength{\emergencystretch}{3em} % prevent overfull lines
\providecommand{\tightlist}{%
  \setlength{\itemsep}{0pt}\setlength{\parskip}{0pt}}
\usepackage{bookmark}
\IfFileExists{xurl.sty}{\usepackage{xurl}}{} % add URL line breaks if available
\urlstyle{same}
\hypersetup{
  hidelinks,
  pdfcreator={LaTeX via pandoc}}

\author{\texorpdfstring{}{}}
\date{}

\begin{document}

\begin{frame}{Structural Finance History: Blake and Alexander Del Mar
Compared}
\protect\phantomsection\label{structural-finance-history-blake-and-alexander-del-mar-compared}
If William Blake provided the visionary, mythopoetic critique of
18th-century monetary metaphysics, the late 19th-century historian
Alexander Del Mar provided its rigorous historiographical counterpart.
Del Mar's \emph{History of Monetary Systems} (1895) operates on a
fundamental thesis that maps onto Blakean cosmology: that the right to
coin money is a sovereign, communal prerogative (a state of Edenic or at
least Beulah-like integration) that has been recurrently hijacked
throughout history by private, atomizing mercantile interests (Urizen)
{[}@delmar1895history{]}.

\begin{block}{The Sacred Sovereign Prerogative of Currency and Coinage
Issuance}
\protect\phantomsection\label{the-sacred-sovereign-prerogative-of-currency-and-coinage-issuance}
For Del Mar, the structural core of monetary history is precisely this
struggle over the prerogative of issuance. When the state controls the
money supply, value is derived from the \emph{societal totality} (the
Human Form Divine) and the collective needs of the people. When private
banks (such as the Bank of England in 1797) or rigid metallic ratios
(Newton in 1717, Liverpool in 1816) assume control, value is transferred
to the dead material of the coin itself or the fictional ledger of the
private creditor {[}@clapham1944bank{]}. Rather than asserting intrinsic
metallic value, Del Mar traces the origins of money to numerary systems
maintained by state authority, a thesis that reveals the transition from
bimetallism to the gold standard not merely as a policy adjustment but
as a metaphysical revolution: the displacement of a communal mechanism
by a rigid, materialist reduction. This perspective illuminates a
\emph{generic} economic principle: every monetary regime encodes a
specific theology of value.
\end{block}

\begin{block}{Blake's Jerusalem Plate 10 and the Architecture of
Enslavement}
\protect\phantomsection\label{blakes-jerusalem-plate-10-and-the-architecture-of-enslavement}
Blake recognized this identical usurpation with prophetic clarity. His
condemnation of the Bank of England, the Mint, and the gold standard was
rooted in his understanding that whoever controls the semiotics of value
controls the human soul. On \emph{Jerusalem} Plate 10, Los refuses to
``Reason \& Compare'' on another man's terms and asserts creative agency
as the ground of valuation (\emph{Jerusalem}, lines 20--21; Erdman
p.~153) {[}@erdman1988complete{]}---the full lines appear in
\hyperlink{sec-four-levels-vision}{§ Fourfold Vision}.

That stance is precisely this recognition that the \emph{architecture}
of representation determines the freedom or subjugation of human life.
The ``System'' Blake must create is not merely a mythological schema but
an entire alternative regime of valuation---one grounded in the
irreducible particularity of creative labor rather than the abstract
universality of the gold sovereign. When Del Mar catalogs the social
consequences of the demonetization of silver---how the contraction of
the monetary base enriched the creditor class while plunging the
producing classes into deflationary poverty---he is documenting the
literal effects of Urizen's ``brazen quadrant'' upon the British and
global economy {[}@delmar1885history{]}.
\end{block}

\begin{block}{The 1844 Bank Charter Act and Ledger-Driven Spectral
Value}
\protect\phantomsection\label{the-1844-bank-charter-act-and-ledger-driven-spectral-value}
The culmination of this architectural subjugation occurred with the 1844
Bank Charter Act (Peel's Act). The fiercely contested debate between the
Banking School and the Currency School narrowed the legal field for note
issue by constraining it under the Bank of England and a statutory
reserve framework with a fixed fiduciary issue. The Currency School
theorists, acting as the ultimate architects of Urizenic geometry,
believed they had finally imprisoned the chaos of economic swerve within
the immutable predictability of the gold standard.

However, by focusing on paper notes, Peel's Act paradoxically ignored
the explosive growth of bank deposits and ledger credit, which later
came to dominate everyday monetary circulation in late 19th-century
Britain. The ``fictive'' or ``spectral value'' against which Blake raged
did not vanish; rather, the ``ghosts'' completely migrated out of the
physical paper artifact and into the invisible, unregulated abstraction
of the bank ledger. This endogenous expansion of invisible credit
represents a profound ontological mutation of the capitalist economy---a
system sustained by mathematical promises that periodically collapsed in
late-century panics.
\end{block}

\begin{block}{The Convergence of Prophetic and Historiographic Critique:
Del Mar and Blake}
\protect\phantomsection\label{the-convergence-of-prophetic-and-historiographic-critique-del-mar-and-blake}
Both Blake and Del Mar understood that the true ``alchemy'' of money had
been stolen from the domain of human creativity. The living, breathing
economy of human interaction had been replaced by a ledger-bound
numerical simulation engineered to perpetually enrich the creditor class
at the direct expense of the producing masses. Reading Blake through
this historiographic lens elevates his poetic epics from obscure
mystical allegories into profound critiques of modern central banking
and financialized alienation {[}@poovey2008genres;
@baucom2005specters{]}. Del Mar's key historiographic insight---that the
demonetization of silver was a deliberate political act designed to
redistribute wealth upward---finds its mythopoetic expression in Blake's
\emph{Jerusalem}, where the architects of financial enclosure are
exposed as agents of Urizenic tyranny.

By the \emph{Milton} injunction against counting gold as ultimate good
(see \hyperlink{sec-illuminated-nonduality}{§ Illuminated Nonduality})
{[}@erdman1988complete{]}, Blake stands alongside Del Mar as a prophet
declaring that the highest sovereign act is not the balancing of a
metallic ratio, but the liberation of the human imagination from the
ledger. Where Del Mar provides the forensic historiography, Blake
provides the visionary horizon: together, they articulate a complete
mythopoetic economics of human liberation.
\end{block}
\end{frame}

\end{document}
